\section{Introduction}

Given any flop $Y_+ \dashrightarrow Y_-$ between two smooth varieties,
it is important to compare their derived categories.
According to a famous conjecture due to Bondal and Orlov~\cite{BO02}, there is expected to be a derived equivalence $\D(Y_+) \simeq \D(Y_-)$.
The aim of this article is to give a new proof of this conjecture for a $7$-dimensional flop using tilting bundles.

\subsection{The Abuaf--Ueda flop}
We~first give the construction of the flop studied in this article.
Consider the $G_2$ Dynkin diagram $\bigcirc\hspace{-0.45em}\waheyheyarrow\hspace{-0.45em}\bigcirc$.
Then by the classification theory of homogeneous varieties, projective homogeneous varieties of the semi-simple algebraic group of type $G_2$
correspond to a marked Dynkin diagram.
The marking $\waheyheytimes\hspace{-0.7em}\waheyheyarrow\hspace{-0.45em}\bigcirc$ corresponds to the $G_2$-Grassmannian
$\G = \Gr_{G_2}$,
the marking $\bigcirc\hspace{-0.45em}\waheyheyarrow\hspace{-0.7em}\waheyheytimes$ corresponds to the $5$-dimensional quadric $\Q \subset \PP^6$,
and the final marking $\waheyheytimes\hspace{-0.7em}\waheyheyarrow\hspace{-0.7em}\waheyheytimes$
corresponds to the (full) flag variety $\F$ of type $G_2$.
There are projections $\F \to \G$ and $\F \to \Q$, and both of~them give $\PP^1$-bundle structures of $\F$.

Now consider the Cox ring of $\F$, namely
\begin{gather*}
C := \bigoplus_{a, b = 0}^{\infty} H^0\big(\F, \stsh_{\F}(a, b)\big) \simeq \bigoplus_{a, b = 0}^{\infty} V_{(a, b)}^{\vee},
\end{gather*}
where $\stsh_{\F}(a, b)$ $\big($resp.~$V_{(a, b)}^{\vee}\big)$ is a line bundle on $\F$ (resp.~the dual of an irreducible representation of $G_2$)
that corresponds to the dominant weight $(a, b)$.
(The bundle $\stsh_{\F}(a, b)$ will be denoted by~$\stsh_{\F}(aH + bh)$ in Section~\ref{geom of homog bdls}.)
Put $C_{a, b} := H^0(\F, \stsh_{\F}(a, b))$ and
\begin{gather*}
C_n := \bigoplus_{a \in \ZZ} C_{n + a, a}.
\end{gather*}
Using these, we can define a $\ZZ$-grading on $C$ by
\begin{gather*}
C = \bigoplus_{n \in \ZZ} C_n.
\end{gather*}
This grading corresponds to the $\Gm$-action on $\Spec C$ obtained from the map $\Gm \to (\Gm)^2$, $\alpha \mapsto \big(\alpha, \alpha^{-1}\big)$
and the natural $(\Gm)^2$-action on $\Spec C$ coming from the original bi-grading.

We~then consider the geometric invariant theory quotients
\begin{gather*}
Y_+ := \Proj(C_+),\qquad
Y_- := \Proj(C_-),\qquad \text{and} \qquad
X := \Spec C_0,
\end{gather*}
where
\begin{gather*}
C_+ := \bigoplus_{n \geq 0} C_n \qquad \text{and} \qquad C_- := \bigoplus_{n \leq 0} C_n.
\end{gather*}
The projective quotients $Y_+$ and $Y_-$ are the total spaces of rank two vector bundles on $\G$ and~$\Q$~respectively.
The affinization morphism $\phi_+\colon Y_+ \to X$ and $\phi_-\colon Y_- \to X$ are small resolutions of the singular affine variety $X$,
which contract the zero-sections.
Furthermore it can be shown that the birational map $Y_+ \dashrightarrow Y_-$ is a $7$-dimensional simple flop with an interesting feature
that the contraction loci are not isomorphic to each other (see~\cite{Ued16}).

The author first learned this interesting flop from Abuaf,
and later noticed that the same flop was independently found by Ueda~\cite{Ued16}.
As such, we call this flop \textit{the Abuaf--Ueda flop}.

Ueda confirmed that Bondal--Orlov conjecture is true for the Abuaf--Ueda flop, using the theory of semi-orthogonal decompositions and their mutations.
However, since there are many other methods to construct an equivalence between derived categories,
it is still an interesting problem to prove the derived equivalence using other methods.

\subsection{Results in this article}
The main purpose of this article is to construct \textit{tilting bundles} on both sides of the flop $Y_+ \dashrightarrow Y_-$,
and construct equivalences between the derived categories of $Y_+$ and $Y_-$ using those tilting bundles.
A tilting bundle $T_{\ast}$ on $Y_{\ast}$ ($\ast \in \{ +, - \}$) is a vector bundle on $Y_{\ast}$ that gives an equivalence
\begin{gather*}
\RHom_{Y_{\ast}}(T_{\ast}, -)\colon \ \D(Y_{\ast}) \to \D(\End_{Y_{\ast}}(T_{\ast}))
\end{gather*}
between two derived categories.
In particular, if we find tilting bundles $T_+$ and $T_-$ with the same endomorphism ring, this then induces an equivalence
$\D(Y_+) \simeq \D(Y_-)$ as desired.

The advantage of this method is that it enables us to study a flop
from the point of view of~the theory of \textit{non-commutative crepant resolutions} (NCCRs)
that were first introduced by~Van den~Bergh~\cite{VdB04b}.
In our case, an NCCR appears as the endomorphism algebra $\End_{Y_{\ast}}(T_{\ast})$ of~a~til\-ting bundle $T_{\ast}$.
Via the theory of NCCRs, we also study the Abuaf--Ueda flop from the moduli-theoretic point of view.

Recall that $Y_+$ and $Y_-$ are the total spaces of rank two vector bundles on $\G$ and $\Q$ respectively.
If there is a variety $Z$ that gives a resolution of an affine variety with rational singularities, and~furthermore
$Z$ is the total space of a vector bundle over a projective variety $W$ that admits a~tilting bundle $T$,
it is natural to hope that the pull back of $T$ via the projection $Z \to W$ gives a~tilting bundle on $Z$.
Indeed, in many known examples, we can produce tilting bundles in such~a~way~\cite{BLV10, H17a, WZ12}.

In our case, it is known that $\G$ and $\Q$ admit tilting bundles (see Section~\ref{geom of homog bdls}).
However, in~our case, the pull-backs of those tilting bundles on $\G$ or $\Q$ do not give tilting bundles on $Y_+$ or $Y_-$.
Thus the situation is different from previous works.
Nevertheless, by modifying bundles that are obtained from tilting bundles on the base $\G$ or $\Q$,
we can construct tilting bundles on $Y_+$ and $Y_-$.
Namely, the tilting bundles we construct are the direct sum of indecomposable bundles that are obtained by taking extensions of
other bundles obtained from $\G$ or $\Q$.
We~check directly that they produce a derived equivalence $\D(Y_+) \simeq \D(Y_-)$.



\subsection{Related works}
If we apply a similar construction to the Dynkin diagrams $A_2$ and $C_2$,
then we obtain the four-dimensional Mukai flop and the (five-dimensional) Abuaf flop~\cite{Seg16} respectively.
Therefore this article shall be viewed a sequel of papers~\cite{H17a, H17b, Seg16}.

Recently, Kanemitsu~\cite{Kan17} classified simple flops of dimension up to eight, which is a certain generalization of a theorem of Li~\cite{Li17}.
His list contains many examples of flops for which the Bondal--Orlov conjecture is still open.
It would be interesting to prove the derived equivalence for all the simple flops that appear in Kanemitsu's list using tilting bundles,
and we can regard this article as a part of such a project.

The Abuaf--Ueda flop is also related to certain (compact) Calabi--Yau threefolds which are studied in~\cite{IMOU16a, IMOU16b, Kuz18}.
Consider the (geometric) vector bundle $Y_+ \to \G$ over $\G$,
then the zero-locus of a regular section of this bundle is a smooth Calabi--Yau threefold~$V_+$ in $\G$.
There is a~similar Calabi--Yau threefold~$V_-$ in~$\Q$.
The papers~\cite{IMOU16b, Kuz18} show that the Calabi--Yau threefolds~$V_+$ and~$V_-$ are L-equivalent, derived equivalent
but are NOT birationally equivalent to~each other.
L-equivalence and non-birationality is due to~\cite{IMOU16b}, and derived equivalence is due to~\cite{Kuz18}.
As explained in~\cite{Ued16}, it is possible to construct a derived equivalence $\D(V_+) \xrightarrow{\sim} \D(V_-)$
for the Calabi--Yau threefolds from a derived equivalence $\D(Y_+) \xrightarrow{\sim} \D(Y_-)$ with a certain nice property.

\subsection{Open questions}
It would be interesting to compare the equivalences in this article and the one constructed by~Ueda.
It is also interesting to find Fourier--Mukai kernels that give equivalences.
In the case of the Mukai flop or the Abuaf flop, the structure sheaf of the fiber product $Y_+ \times_X Y_-$ over the singularity $X$
gives a Fourier--Mukai kernel of an equivalence (see~\cite{H17b, Ka02, Na03}).
Thus it is interesting to ask whether this fact remains to hold or not for the Abuaf--Ueda flop.

Another interesting topic is to study the autoequivalence group of the derived category.
Since we produce some derived equivalences that are different to each other in this article, we can find some non-trivial autoequivalences
by combining them.
It would be interesting to find an action of an reasonable group on the derived category of $Y_+$ (and $Y_-$)
that contains our autoequivalences.

